%! Author = lili
%! Date = 9/7/26


\section{Nachweis eines Reportergenproduktes in transgenen Pflanzen}\label{sec:nachweis-eines-reportergenproduktes-in-transgenen-pflanzen}

\subsection{Einführung}\label{subsec:einfuhrung}

\defin{genexpressionsanalyse}

\ausschnitt{Die Genexpressionsanalyse kann sowohl für einzelne Transkripte als auch für das ganze
Transkriptom angewendet werden und ermöglicht qualitative und
quantitative Aussagen über die Aktivität der Gene.
Im Kandidatengenansatz muss nur die Sequenzinformation des zu
untersuchenden Transkripts bekannt sein.
Für die Transkriptom-Analyse wird die gesamte Sequenzinformation des
Transkriptoms gebraucht (oder gleichzeitig generiert).}{folien21}
